\input{../tex-inputs/latex-leseansicht-vorspann}

\section[Arthur Schnitzler an Gustav Schwarzkopf, {[}16. oder 17{]}. 7. 1906]{L04377 Arthur Schnitzler an Gustav Schwarzkopf, {[}16. oder 17{]}. 7. 1906}
\nopagebreak\mylabel{L04377v}
\rehead{ }\normalsize\beginnumbering\briefempfaengerindex{Schwarzkopf, Gustav@\textsc{Schwarzkopf, Gustav}!zzzSchnitzler, Arthur@\emph{von Arthur Schnitzler}!1906-07-172@{{[}16. oder 17{]}. 7. 1906}|(be}
\toendnotes[C]{\smallbreak\pagebreak[2]}
\correspDesc{Versand  durch Arthur Schnitzler im Zeitraum [16. oder 17]. 7. 1906 in Kurhotellet
\newline{}Übermittlung  am 17. 7. 1906 in Helsingør
\newline{}Erhalt  durch Gustav Schwarzkopf im Zeitraum [18. 7. 1906 – 22. 7. 1906?] in Tiefer Graben 23}\toendnotes[C]{\smallbreak}
\Standort{DLA, A:Schnitzler, HS.1985.1.1897.}
\physDesc{Bildpostkarte, 447 Zeichen
\newline{}Handschrift: Bleistift, deutsche Kurrent
\newline{}Versand: Stempel: »\nobreak{}\oindex{Helsingør@\textbf{Helsingør}|pwk}Helsingør, 17. 7. 1906, 2–5 E\nobreak{}«.  }\toendnotes[C]{\smallbreak}\pstart{}{\pb}\textsc{Austria}\oindex{Österreich@\textbf{Österreich}|pw}\pend{}\pstart{}Hrn \textsc{Gustav Schwarzkopf}\pend{}\pstart{}\textsc{Wien} I\oindex{Wien@\textbf{Wien}|pw}\pend{}\pstart{}Tiefer Graben 23\oindex{Wien@\textbf{Wien}!I., Innere Stadt@\textbf{I., Innere Stadt}!Tiefer Graben 23@\textbf{Tiefer Graben 23}|pw}.\pend{}{\bigskip}
\pstart
           \noindent{}\centering{}{\pb}\textcolor{gray}{\textbf{Hilsen fra Helsingør\oindex{Helsingør@\textbf{Helsingør}|pw}}}\hspace*{2.5em}\textcolor{gray}{\textbf{Salut for Kronborg\oindex{Schloss Kronborg@\textbf{Schloss Kronborg}|pw}}}\pend
           \vspace{1em}
\pstart
           \noindent{}{\pb}lieber Guſtav, auch hier hat ſich das Wetter zum ſchlechten gewandt,
               aber die Annehmlichkeiten des Hotels\oindex{Kurhotellet@\textbf{Kurhotellet}|pw} laſſen es ertragen\pend
           
\pstart
           Von
               Paul M.\pwindex{Marx, Paul 21.\,7.\,1879 Wien – 30.\,10.\,1956 ebd.@\textsc{Marx, Paul} (21.\,7.\,1879 Wien – 30.\,10.\,1956 ebd.)|pw} kein Sterbenswort an uns{ }ſeit wir fort{ }ſind.
                  Ich arbeite, und zwar richtig ein Stück\pwindex{Schnitzler, Arthur 15. 5. 1862 Wien – 21. 10. 1931 ebd.@\textsc{Schnitzler, Arthur} (15. 5. 1862 Wien – 21. 10. 1931 ebd.)!Wort. Tragikomödie in fünf Akten@\strich\emph{Das Wort. Tragikomödie in fünf Akten}|pwv},
            wie Sie vorausgeſagt haben. Wahrhaftig, man möchte nicht alle Ihre Prophezeihungen ke{\geminationn}en.\pend
           
\pstart
           {\pb}\label{K_L04377-1v}\edtext{Heini\pwindex{Schnitzler, Heinrich 9.\,8.\,1902 Hinterbrühl – 12.\,7.\,1982 Wien@\textsc{Schnitzler, Heinrich} (9.\,8.\,1902 Hinterbrühl – 12.\,7.\,1982 Wien)|pw} zeichnet}{\lemma{\textnormal{\emph{Heini zeichnet}}}\Cendnote{\textnormal{Vgl. A. S.: \emph{Tagebuch}, 14. 7. 1906.}}}\label{K_L04377-1}: Eiſenbahnen,
               \textcolor{gray}{Zeiger} und Eſel.\pend
           \pstart Herzliche Grüße Ihnen u Dr Max\pwindex{Schwarzkopf, Max 12.\,6.\,1857 Wien – 14.\,4.\,1928 ebd.@\textsc{Schwarzkopf, Max} (12.\,6.\,1857 Wien – 14.\,4.\,1928 ebd.)|pw}{ }\textcolor{gray}{von uns beiden.}{ }\spacefill\mbox{A.}\pend{}\selectlanguage{ngerman}\endnumbering\briefempfaengerindex{Schwarzkopf, Gustav@\textsc{Schwarzkopf, Gustav}!zzzSchnitzler, Arthur@\emph{von Arthur Schnitzler}!1906-07-162@{{[}16. oder 17{]}. 7. 1906}|)be}\mylabel{L04377h}
\begin{anhang}
\end{anhang}\newcommand{\dateiname}{L04377}\newcommand{\titel}{Arthur Schnitzler an Gustav Schwarzkopf, [16. oder 17]. 7. 1906}\newcommand{\editorInnen}{Herausgegeben von Jahnke, SelmaMüller, Martin Anton}\input{../tex-inputs/latex-leseansicht-abspann}
